\section{Experimental Evaluation}\label{sec:experiments}
Single-threaded C++17 measurements use \texttt{g++ 14.2 -O3} on an
AMD EPYC 9V74 shared container. Exact/KB times are medians of seven runs;
hybrid times and outputs average 100 runs, including exact returns.
Graph construction and verification are untimed. Source code, graph generators, seeds, edge lists, and full results
for the comparison are available at
\url{https://github.com/aarashahadiacademic-1/tanner-cycle-counting}.

\subsection{Benchmark Construction for Comparison with KB
\cite{karimi2013}}
The graphs in Table~\ref{tab:comparison-exp} are synthetic sparse parity-check
matrices, not standardized codes or decoding benchmarks.
We favor small degrees to limit the branching of our local search.

The labels in the first column of Table~\ref{tab:comparison-exp}
identify the graph construction and, where applicable, its lift.
Bases A--C use point--line incidence $y=sx+b$ over $\mathbb F_Q$,
the finite field with $Q$ elements, with
$0\le x \le d_c-1$ and $0\le s\le d_v-1$.
Their parameters $(Q,d_v,d_c)$ are $(7,3,6)$, $(5,3,4)$, and
$(11,4,8)$, respectively. This construction gives the stated degrees
and excludes $4$-cycles.
Base D is obtained from the symplectic generalized quadrangle over
$\mathbb F_5$ by splitting each degree-six point into two degree-three
variables. It has 312 variables, 156 checks, degree pair $(3,6)$,
and girth eight.

For A--D, labels specify the base, lift size, and cyclic ($c$) or
random ($r$) lift.
R1200 uses random socket pairing followed by
degree-preserving switches to remove $4$-cycles.
P1200 and P1600 use degree-constrained progressive edge growth.
The numbers in these R/P labels denote the number of variables.

Q410 and Q412 are $(3,4)$ quasi-cyclic graphs with girths 10 and 12
and circulant sizes 10,112 and 9,088, respectively.
They use the first four columns of Zhang and Wang's exponent
matrix~\cite{zhang2010}, with modified circulant shifts.

Construction details, exponent matrices, and fixed seeds are provided
in the repository linked above. Every realized graph was checked for
simplicity, connectivity, degrees, and girth.

Our KB implementation follows Karimi and
Banihashemi~\cite{karimi2013}, rather than their
software, using exponent-vector messages,
node sums, reverse subtraction, nonzero-message
processing, root retransmission suppression,
and earlier-root deactivation.
Table~\ref{tab:comparison-exp} compares the three methods on $N_g$.
Separate longer-target comparisons in the repository show complementary
performance: KB can be faster on some longer cycle-rich targets.
These implementation-dependent results on synthetic graphs
do not establish superiority across LDPC families.

\subsection{Hybrid accuracy and runtime}
Table~\ref{tab:comparison-exp} compares all three algorithms on the
same target $N_g$, using $\varepsilon=0.08$, $r=101$, and a fixed
interleaving quantum $Q=256$. Here $n$ is the number of variable
nodes and $N_g$ is the exact girth-cycle count.
KB and Exact times are medians of seven runs, while Hybrid times
and outputs average 100 runs. Hybrid time includes both interleaved
procedures and scheduling overhead. MARE is the mean of
$100|\widetilde N_g-N_g|/N_g$ over all 100 outputs,
including exact returns.

\begin{table*}[t]
\centering
\caption{Same-target comparison for $N_g$; running times in milliseconds.
KB and Exact are medians of seven runs; Hybrid and MARE average 100 runs.
MARE is mean absolute relative error in percent.}
\label{tab:comparison-exp}
\footnotesize
\setlength{\tabcolsep}{5pt}
\begin{tabular}{lccrrrrrr}
\hline
Graph & $(d_v,d_c)$ & $g$ & $n$ & $N_g$
& KB & Exact & Hybrid & MARE (\%)\\
\hline
A16c & $(3,6)$ & 6 & $672$ & $2,560$ & 1.507 & 0.497 & 0.799 & 7.7\\
A128r & $(3,6)$ & 6 & $5,376$ & $176$ & 58.701 & 6.215 & 104.750 & 0.0\\
A128c & $(3,6)$ & 6 & $5,376$ & $20,480$ & 14.842 & 5.699 & 1.382 & 8.8\\
B128c & $(3,4)$ & 6 & $2,560$ & $4,608$ & 3.068 & 1.095 & 0.454 & 7.6\\
C128c & $(4,8)$ & 6 & $11,264$ & $204,928$ & 112.565 & 44.423 & 3.196 & 7.6\\
D16c & $(3,6)$ & 8 & $4,992$ & $57,680$ & 75.800 & 18.504 & 4.548 & 6.9\\
R1200 & $(4,8)$ & 6 & $1,200$ & $1,564$ & 47.418 & 4.171 & 24.667 & 6.5\\
P1200 & $(3,6)$ & 8 & $1,200$ & $65$ & 44.640 & 4.211 & 44.721 & 0.0\\
P1600 & $(3,4)$ & 10 & $1,600$ & $11$ & 70.511 & 5.195 & 80.908 & 0.0\\
Q410 & $(3,4)$ & 10 & $40,448$ & $40,448$ & 1125.038 & 128.667 & 90.153 & 7.5\\
Q412 & $(3,4)$ & 12 & $36,352$ & $308,992$ & 2971.170 & 292.054 & 62.125 & 7.6\\
\hline
\end{tabular}
\end{table*}

Zero MARE indicates that exact counting finished first in every run,
as on A128r, P1200, and P1600. Cycle-rich instances can benefit
from sampling, whereas sparse ones may incur interleaving overhead
before exact termination. On R1200, sampling wins all 100 runs,
but the hybrid is slower than standalone exact counting because
state-machine steps have unequal practical costs. At girths 10
and 12, Q410 and Q412 use inverse sampling in every run and
obtain lower hybrid latency than either exact method. These
implementation-dependent results do not establish superiority
across LDPC families.

The observed MARE values around $8\%$ are consistent with a
normal approximation to uncapped inverse sampling. Since $S_r$
has mean $r/p$ and variance $r(1-p)/p^2$, for large $r$ the
relative estimation error is approximately normal with variance
$(1-p)/r$. Hence the expected absolute relative error is
approximately $\sqrt{2(1-p)/(\pi r)}$, or about $7.94\%$ for
$r=101$ and small $p$. This is an asymptotic approximation for
uncapped inverse sampling, not an exact hybrid guarantee;
exact termination can reduce the overall MARE. The guarantee
in Theorem~\ref{thm:hybrid} concerns expected error, so a
100-run empirical MARE may exceed $\varepsilon$.
